\section{Das atomare Modell}

Auf dem normierten Grundintervall $u\in[0,1]$ sei
\[
  T(u)=1-2u.
\]
Damit fällt $T$ gleichförmig von $1$ nach $-1$. Die unmittelbar folgenden
Eigenschaften sind
\[
  T(0)=1,
  \qquad
  T(1)=-1,
  \qquad
  T(u)^2\leq 1.
\]
Durch Spiegelung des Grundintervalls entsteht eine periodische Dreieckswelle.
Für einen Winkel $\theta\in[0,\pi]$ entspricht die Normierung
$u=\theta/\pi$ der Darstellung
\[
  T\!\left(\frac{\theta}{\pi}\right)
  =1-\frac{2\theta}{\pi}.
\]

Diese Steuerung ist horizontal linear. Wird sie unmittelbar auf eine Ellipse
gehoben, ist die vertikale Bewegung dagegen an den Umkehrpunkten nicht regulär:
Die zugehörige vertikale Ableitung wird unbeschränkt. Genau dieser Kollaps
motiviert den später eingeführten Glättungsparameter.
